% ==============================================================================
% SECCIÓN 9: CAUCES DE SEGUIMIENTO CON EL CLIENTE
% ==============================================================================
\newpage
\section{Cauces de Seguimiento con el Cliente}


\begin{center}
\small
\begin{tabularx}{\textwidth}{|>{\raggedright\arraybackslash}p{3cm}|>{\raggedright\arraybackslash}p{2.6cm}|>{\raggedright\arraybackslash}X|>{\raggedright\arraybackslash}p{3cm}|}
\hline
\rowcolor{tableHeader}
Cauce propuesto & Frecuencia & Información / propósito & Participación del cliente \\
\hline
Resumen de estado por correo electrónico & Semanal & Trabajo completado, siguiente paso, bloqueos, riesgos y decisiones necesarias. & Nombrar contacto destinatario y responder a decisiones solicitadas. \\
\hline
Reunión de seguimiento & Cada dos semanas & Revisar avances, mostrar el prototipo o versión disponible, resolver dudas y acordar cambios. & Asistir con el interlocutor designado y validar acuerdos. \\
\hline
Registro compartido de decisiones e incidencias & Actualización continua & Mantener identificador, descripción, impacto, responsable y estado de cada decisión o incidencia. & Revisar y confirmar las decisiones que requieran aprobación del centro. \\
\hline
Reunión de inicio con el centro (realizada el 25/09/2026, en la ETSIIT) & Una vez, al inicio & Exposición en clase sobre el centro, el alumnado PTVAL y la necesidad que origina la app; sesión de preguntas y respuestas con el alumnado de la asignatura. Origen de los requisitos iniciales. & El tutor y la profesora del centro presentaron la necesidad y resolvieron las dudas planteadas. \\
\hline
Visita al centro (realizada el 02/10/2026 por nuestro compañero de grupo) & Una visita realizada & Conocer el contexto real: espacios, dinámica de las clases, apoyos que ya utilizan y forma de trabajar con el alumnado. Contrastar los requisitos con la realidad del centro. & Recibir al miembro del equipo y mostrar el centro y su funcionamiento. \\
\hline
Material del centro en PRADO & Consulta continua & Documentación aportada por el centro como referencia para requisitos y apoyos visuales. & Aportar material adicional si lo considera útil. \\
\hline
Orientación del profesorado de DGP y MDA & Sesiones de clase y tutorías & Orientación sobre el enfoque del proyecto y la documentación. No sustituye las decisiones del centro. & Ninguna (cauce interno académico). \\
\hline
\end{tabularx}
\end{center}

El equipo enviará un aviso extraordinario cuando una incidencia pueda afectar a un requisito aprobado, una fecha objetivo o la seguridad de los datos. Se registrará la decisión del centro en el acta o registro correspondiente; los cambios de alcance no se darán por aprobados solo por comunicarlos en el informe de estado.